% =============================================================================
%  CHAPTER 2 — BACKGROUND & RELATED WORK
%
%  >>> AI-ASSISTED FIRST DRAFT (generated from YOUR docs/RELATED_WORK.md). <<<
%  This is raw material, NOT submittable text. Before submission you MUST:
%    1. Rewrite it in your own voice (the markers quiz you on this in the viva).
%    2. Verify every citation against the actual paper — several entries in
%       references.bib are marked `VERIFY` (authors, venues, the SWE-Refactor ID
%       and its 41.6% figure, on which your motivation depends).
%    3. Disclose AI drafting assistance in the Declarations page (main.tex).
%  RUBRIC: Framing of Research Problem (15%). BUDGET: ~8-12 pages.
% =============================================================================
\chapter{Background and Related Work}
\label{ch:background}

\section{Technical Background}
\label{sec:bg-technical}

This section introduces the concepts a reader needs to follow the rest of the
report: the structural representation an IDE maintains for source code, the
refactoring engine that operates on it, and the tool-use mechanism by which a
large language model (LLM) can invoke external operations.

\paragraph{Program structure and the PSI.} Modern IDEs do not treat source code
as text. IntelliJ IDEA maintains a \emph{Program Structure Interface} (PSI): a
typed abstract syntax tree augmented with a project-wide reference index that
resolves each identifier to its declaration and each declaration to all of its
usages. Because the PSI encodes types, scopes, and references rather than
character offsets, an operation expressed over it can reason about \emph{which}
\texttt{addVisit} a call refers to, or \emph{which} overload of \texttt{getPet}
a signature selects --- information that is absent from the raw text.

\paragraph{Refactoring engines.} A refactoring is a behaviour-preserving
transformation of a program. IntelliJ implements each refactoring as a
\emph{processor} (e.g.\ \texttt{RenameProcessor},
\texttt{MoveClassesOrPackagesProcessor}, \texttt{ChangeSignatureProcessor})
that mutates the PSI and propagates the change to every dependent reference,
import, and --- where requested --- comment and string occurrence, atomically.
The correctness of the result follows from the engine's semantic model, not from
the caller. This is the foundation of the design: by constraining the agent to \emph{only}
drive these PSI-backed operations --- never to edit source text directly ---
correctness no longer depends on the agent's own reasoning about where a change
must ripple. The agent need only choose \emph{which} refactoring to apply; the
engine then guarantees that the rename, signature change, or move is
behaviour-preserving by definition, propagating to every dependent reference
across the codebase atomically and without the missed call sites, broken imports,
or other unintended side effects that ad hoc text editing risks.

\paragraph{LLM tool use and the Model Context Protocol.} Contemporary LLMs
support \emph{function calling}: given a set of named tools described by a JSON
schema, the model emits a structured call (a tool name and a JSON argument
object) that the host application executes, returning the result for the model
to continue reasoning. The Model Context Protocol (MCP) standardises this
exchange so that any compliant agent can discover and call a server's tools.
This project exposes IntelliJ's refactoring engine through exactly such a
tool interface.

\section{LLM-Driven Code Refactoring}
\label{sec:bg-llm-refactoring}

A growing body of work evaluates how well LLMs can be applied to software
engineering and code tasks; \citet{hou2024llm4se} provide a systematic literature
review of 395 such studies and identify interpretability and trustworthiness,
together with limited evaluation on industrial data, as open challenges that still
impede real-world adoption --- challenges that a structurally-safe,
behaviour-preserving execution layer is well placed to lower.
\citet{cordeiro2024refactoring} study StarCoder2 across thirty open-source Java
projects, and find that the model outperforms developers on systematic, repetitive
refactorings but falls short on context-heavy, architectural ones --- precisely the
cases that depend on project-wide reference information. \citet{potentialllm2024} examine 180
real-world refactorings across twenty projects and report that ChatGPT
identifies only 28 of them unaided, yet reaches an 86.7\% success rate once the
refactoring sub-category is pre-specified. This last result is important for the
present work: constraining the model to a fixed menu of operation types --- as a
typed tool API does --- substantially improves reliability over free-form
generation.

The clearest statement of the problem comes from repository-level benchmarks.
\citet{swerefactor2025} assemble 1{,}099 developer-written, behaviour-preserving
Java refactorings validated by compilation, tests, and RefactoringMiner, and
report a best LLM pass rate of only \textbf{41.6\%}. This ceiling on
text-generation approaches is the empirical motivation for delegating execution
to an IDE. Complementarily, \citet{llmrefactorlimits2025} catalogue the failure
modes behind that ceiling --- hallucinated edits, broken semantics, and
\emph{scope blindness} (missing usages in other files) --- and identify IDE
integration as the natural mitigation. This work targets the \emph{execution} step those failure modes share: it removes
free-text generation from the write path entirely, so the scope-blind,
semantically-unsafe edits behind that ceiling --- missed cross-file usages, broken
imports, wrong-overload renames --- cannot arise, because every change is applied by
an AST-aware engine that propagates it across the whole codebase by construction.

\section{Program Transformation Agents}
\label{sec:bg-agents}

A parallel line of work embeds LLMs in autonomous agents that act on a codebase
through a tool interface. \citet{sweagent2024} introduce the Agent--Computer
Interface (ACI): a small, shell-like toolset (view, search, edit, run) through
which an agent resolves real GitHub issues, achieving 12.5\% on SWE-bench at
release. The ACI is the closest architectural analogue to this project's HTTP
tool surface, but its \texttt{edit} tool performs raw text replacement with no
awareness of program structure. \citet{openhands2025} generalise this to a
sandboxed platform in which agents drive a containerised operating system
(shell, browser, Python REPL); the action space is maximally general but
correspondingly unconstrained. \citet{autocoderover2024} narrow the focus to
program repair, using AST-backed code-search APIs (class signatures, method
bodies, call graphs) to locate context before generating a patch; notably,
however, the AST is used only for \emph{read-only retrieval}, while the patch is
still applied as a text diff. Across this line of work the pattern is
consistent: structure, where used at all, informs \emph{reading}, never
\emph{writing}.

% --- AI FIRST DRAFT (from arXiv:2511.04824 + your research notes). Verify the
%     figures against the paper and edit into your voice. ---
The most direct empirical evidence for how such agents refactor in practice comes
from \citet{agenticrefactoring2025}, who conduct the first large-scale study of
production coding agents (Codex, Cursor, Devin, Claude Code), detecting 15{,}451
refactorings across 14{,}988 real Java commits with RefactoringMiner and assessing
their structural impact with DesigniteJava. Two of their findings directly motivate
the present work. First, although agents refactor in 26.1\% of commits, they
overwhelmingly perform \emph{localized, low-level} edits --- change-variable-type
(11.8\%), rename-parameter (10.4\%), and rename-variable (8.5\%) dominate --- a
``preference for localized improvements over the high-level design changes common
in human refactoring''. The structural operations agents largely avoid (move,
extract, signature change) are precisely those where free-text editing is least
safe, as the present evaluation confirms (\Cref{ch:eval}); a structured execution
layer is what makes them safe to attempt. Second, their quality analysis finds only
\emph{small} structural gains (e.g.\ a median reduction of 15.25 class lines of
code) and reports that agents ``fail to meaningfully reduce actual design
smells''; moreover 53.9\% of the detected refactorings occur in \emph{tangled}
commits with no explicit refactoring intent --- applied incidentally alongside
other edits rather than as deliberate, isolated transformations. The study does
not establish that these edits compile or preserve behaviour. These are precisely
the guarantees an IDE-backed execution layer provides by construction, and the
intentful, single-purpose nature of a typed refactoring operation is what makes
each change individually reviewable rather than entangled with unrelated edits.
This project therefore complements their descriptive account of \emph{what} agents
refactor with a prescriptive question: does constraining the agent's write action
to structured, semantics-preserving operations improve correctness, reproducibility,
and auditability?

\section{AST-Safe and Structured Refactoring with AI}
\label{sec:bg-structured}

The systems most closely related to this project combine an LLM with an IDE's
refactoring engine so that execution is structurally safe. \citet{emassist2024}
present EM-Assist, an IntelliJ plugin in which an LLM suggests code ranges to
extract and IntelliJ's PSI engine validates and performs the extraction,
reaching 53.4\% Recall@5 on 1{,}752 real refactorings with a 94.4\% positive
developer rating. EM-Assist is the most architecturally similar prior system and
the natural baseline for comparison; its limitation is that it is
\emph{single-operation} (extract-method only) and hardwired inside the IDE, with
no external interface an agent could call.
\citet{movemethod2025} combine LLM suggestions with IDE static analysis and
retrieval over prior refactorings to recommend Move Method, reaching 73\%
Recall@1 --- again a single operation, exposed through no agent-callable API.
\citet{mantra2025} orchestrate multiple LLM agents with contextual RAG and
self-repair across six method-level refactoring types (including extract-method),
with 82.8\% of generated refactorings compiling and passing tests overall; the
pipeline is sophisticated but still produces text-level replacements rather than
driving the IDE's engine.
A distinct, human-centred strand examines how developers and an LLM actually
\emph{converse} to decide and communicate which refactoring to apply.
\citet{refactoringconversations2024} text-mine 17{,}913 developer--ChatGPT prompts
and responses about refactoring, finding that developers typically issue generic,
natural-language requests while the model supplies the specific refactoring
intention, and that requests range from simply pasting code (2.3\%) to detailed
descriptions with explicit fix instructions (41.9\%). That work and the present one
are complementary: it concerns \emph{how} the human and the model jointly decide and
communicate which refactoring to propose, whereas this project concerns \emph{how
the chosen refactoring should be executed} so that the result is correct by
construction regardless of which party --- human or model --- selected it.
The gap this work addresses follows directly. To the best of our knowledge, no prior
system exposes the IDE's PSI-backed, AST-aware refactoring engine as a
\emph{multi-operation}, \emph{externally-callable} tool API through which an arbitrary
coding agent --- Claude Code, or any other function-calling LLM --- can \emph{execute}
correct, project-wide refactorings, rather than merely read the AST for context before
emitting a text patch. The prior systems surveyed here are single-operation, hardwired
inside the IDE, or text-level; the closest contemporaneous work is JetBrains' own MCP
server, which exposes IDE capabilities to MCP clients. The present system is
complementary to that effort, contributing a model-agnostic, multi-operation
refactoring tool surface of exactly this kind (\Cref{ch:design}).

\section{Constrained Action Spaces and Tool Use}
\label{sec:bg-action-space}

\citet{demystifying2025} question whether elaborate autonomous agents are needed at
all: their \emph{Agentless} approach shows that a deliberately simple, constrained
pipeline --- localise, repair, validate --- outperforms more complex agent designs,
and they observe that over-rich tool and action spaces lead models into imprecise or
incorrect tool use. This evidences the central tension between \emph{generality} (a
Turing-complete shell) and \emph{safety} (a constrained action space), and positions
the present work: rather than exposing a general shell, the plugin offers only typed,
semantically defined refactoring operations, trading generality for guaranteed
structural validity. At the other end of the spectrum, the widely used Aider
assistant~\citep{aider} edits whole files or diffs committed to git with no AST
awareness, exemplifying the state-of-practice text layer that this project
contrasts with the semantic (symbol) layer.

\section{Refactoring Detection as an Evaluation Oracle}
\label{sec:bg-detection}

Finally, the evaluation in Chapter~\ref{ch:eval} draws on refactoring
\emph{detection}. \citet{refactoringminer2025} present RefactoringMiner, the
state-of-the-art tool for detecting refactoring operations in commit histories
via refactoring-aware AST differencing; it serves as a ground-truth oracle in
most of the benchmarks above and is used here to confirm that a tool call
produced the intended refactoring type. \citet{motivations2025} complement this
with an LLM-driven taxonomy of \emph{why} developers refactor, which, while not
central to this work, suggests when an agent might appropriately propose a
refactoring. Compilation --- the primary oracle in this work --- is a necessary
but not sufficient check on behaviour preservation. A stronger, complementary
signal is functional equivalence: \citet{difffuzzing2026} apply differential
fuzzing to LLM-generated refactorings, executing the original and refactored code
on shared inputs to surface behavioural divergence that compilation alone cannot
detect. Such behavioural oracles are orthogonal to the structural guarantees
studied here --- the structured engine guarantees a semantics-preserving
\emph{transformation}, whereas differential fuzzing would \emph{independently
verify} that preservation --- and combining the two is a natural direction for
future evaluation (\Cref{ch:conclusion}).

\section{Summary and Research Gap}
\label{sec:bg-gap}

Table~\ref{tab:gap} situates this project against the prior work surveyed above.
Text-generation approaches offer broad operation coverage but no correctness
guarantee and a measured ceiling near 41.6\%~\citep{swerefactor2025};
single-operation IDE plugins such as EM-Assist~\citep{emassist2024} guarantee
correctness but expose one operation and no external API; and agent platforms
such as SWE-agent~\citep{sweagent2024} provide an external tool interface but
edit as text.

\begin{table}[h]
\centering
\caption{Where this project sits relative to prior approaches.}
\label{tab:gap}
\begin{tabular}{@{}llll@{}}
\toprule
Approach & Example & Execution & Gap addressed here \\
\midrule
LLM generates code text   & SWE-agent, Aider, MANTRA & text diff      & no correctness guarantee \\
Single-op IDE plugin      & EM-Assist, Move Method   & IntelliJ PSI    & one operation; no external API \\
Read-only AST + patch     & AutoCodeRover            & text patch      & AST not used for writes \\
\textbf{This project}     & ---                      & IntelliJ PSI    & multi-op, externally callable \\
\bottomrule
\end{tabular}
\end{table}

The pattern across these rows is consistent: the systems that \emph{execute} safely
cannot be called by an agent, and the systems an agent \emph{can} call execute as
text. Even the approaches that exploit the abstract syntax tree use it only to
\emph{read} --- to gather context before emitting a textual patch --- never to
\emph{write}. This project closes that gap. It exposes the IDE's PSI engine as a
multi-operation, externally-callable tool surface, so that a coding agent does not
merely read the AST but \emph{performs its edits through it}: every rename, move, or
signature change is applied by the engine and propagated across the whole codebase,
making the result correct by construction. To the best of our knowledge, allowing an
external agent to drive a general, multi-operation IDE refactoring engine in this way
has not been done before; it is the foundation of the design (\Cref{ch:design}) and
evaluation (\Cref{ch:eval}) that follow.
